\chapter{ЧИСЛЕННОЕ ИНТЕГРИРОВАНИЕ ДИФФЕРЕНЦИАЛЬНЫХ УРАВНЕНИЙ С ПОМОЩЬЮ GPU}
\label{cha:gpu}

\section{Постановка задачи}

Центральная идея проекта --- перенос расчёта физической динамики на графический
ускоритель. Графические процессоры содержат тысячи вычислительных ядер и при
правильной организации данных способны на порядки ускорять массовые однотипные
вычисления, к которым относится интегрирование уравнений движения множества тел.
Задача данной работы --- реализовать решатель (солвер) дифференциальных уравнений
движения \eqref{eq:newton-euler} на GPU, функционально эквивалентный эталонной
CPU-реализации (раздел~\ref{cha:cpu}). Работа выполнена в репозитории
\texttt{dyphur}.

\section{Переносимый вычислительный слой}

Чтобы не привязываться к одному поставщику оборудования, GPU-вычисления
организованы поверх стандарта SYCL~2020~\cite{sycl2020} (реализация
AdaptiveCpp~\cite{alpay2020adaptivecpp}). Это позволяет
из единого исходного кода на C++ получать исполняемый код для нескольких
бэкендов:
\begin{itemize}
  \item \textbf{CUDA} --- ускорители NVIDIA;
  \item \textbf{HIP} --- ускорители AMD (ROCm);
  \item \textbf{OpenMP} --- исполнение на CPU (эталон, раздел~\ref{cha:cpu}).
\end{itemize}

Над SYCL построен тонкий слой абстракции (модуль \texttt{compute}),
предоставляющий: владение устройством и очередью (\texttt{Device}),
типизированные буферы памяти с загрузкой/выгрузкой (\texttt{Buffer<T>}),
параллельный запуск ядер (\texttt{parallel\_for}), а также детерминированные
коллективные примитивы --- редукцию (\texttt{reduce}), сортировку по ключу
(\texttt{sort\_by\_key}) и последовательно-согласованное атомарное сложение
(\texttt{atomic\_add\_seq}). В код устройства не проникают контейнеры STL,
исключения и виртуальные функции; контроль обеспечивается отдельной проверкой
зависимостей на этапе сборки. Линейная алгебра в ядрах выполнена на
SoA-совместимых типах (\texttt{Vec3f}, \texttt{Quatf}, \texttt{Mat3f},
\texttt{Transform}).

\section{GPU-интегратор}

Интегрирование уравнений движения распараллелено по телам: каждому телу
сопоставляется отдельный рабочий элемент (work-item), который выполняет
симплектический шаг \eqref{eq:symplectic} и обновление кватерниона ориентации.
Раскладка SoA обеспечивает коалесцированный доступ к глобальной памяти ---
соседние рабочие элементы обращаются к соседним ячейкам массивов координат и
скоростей.

\begin{lstlisting}[language=C++,caption={Ядро интегрирования скоростей и позиций на GPU (схематично)},label=lst:gpu]
void integrate_gpu(Stream& s, BodyView b, float h, int n_substeps) {
    const size_t n = b.count;
    for (int sub = 0; sub < n_substeps; ++sub) {
        parallel_for(s, n, [=](size_t i) {
            if (b.inv_mass[i] == 0.0f) return;       // статические тела пропускаем
            // v += h * (g + M^-1 F)   — симплектический Эйлер
            b.vel_x[i] += h * b.gx;  b.vel_y[i] += h * b.gy;  b.vel_z[i] += h * b.gz;
            // x += h * v
            b.pos_x[i] += h * b.vel_x[i];
            b.pos_y[i] += h * b.vel_y[i];
            b.pos_z[i] += h * b.vel_z[i];
            // q += h * 0.5 * (omega ⊗ q), затем нормировка
            integrate_orientation(b, i, h);
        });
    }
}
\end{lstlisting}

\section{Детерминизм на GPU}

Параллельные вычисления с плавающей точкой не ассоциативны: порядок суммирования
влияет на результат. Чтобы обеспечить побитовую воспроизводимость, в коллективных
операциях применяются детерминированные схемы: редукция выполняется по
фиксированному дереву свёртки, сортировка --- битоническая (детерминированный
порядок сравнений), счётчики --- через последовательно-согласованное атомарное
сложение. В результате на фиксированной паре <<сборка\,--\,оборудование>> GPU-ядро
даёт одинаковый хеш состояния при повторных запусках (см.
раздел~\ref{cha:test}).

\screenshot{0.85}{Сравнение CPU- и GPU-интегратора}{Снимок экрана: сопоставление
траекторий/состояний эталонной CPU-симуляции и GPU-симуляции на одинаковом
сценарии (совпадение поведения, воспроизводимость), либо профиль загрузки
GPU.}{gpu-compare}

\section{Результаты}

Реализован решатель дифференциальных уравнений движения на GPU поверх
переносимого вычислительного слоя SYCL/AdaptiveCpp с бэкендами CUDA и HIP и
CPU-фолбэком. GPU-интегратор функционально эквивалентен эталонной
CPU-реализации, обеспечивает детерминизм результата и в составе эталонного
сценария (раздел~\ref{cha:test}) выполняет симуляцию 1024~тел в реальном времени.
Создан задел для масштабирования числа тел и полного переноса конвейера
обнаружения столкновений и решателя ограничений на GPU на последующих этапах.
